% AA-14 | Algebra abstrakcyjna | Iteracja VII
\chapter[Pierścienie artinowskie i lokalne]{Pierścienie artinowskie i lokalne oraz lemat Nakayamy}
\label{aa:ch:artin-local}
\setlength{\parskip}{2pt}

\section{Dwa przeciwne warunki łańcucha}
Pierścienie noetherowskie charakteryzują się stabilizacją
\emph{wstępujących} ciągów ideałów. Można też wymagać,
aby stabilizowały się ciągi zstępujące. Te dwa warunki
nie są równoważne, a ich konsekwencje są bardzo różne.
W całym rozdziale pierścienie są przemienne, mają jedynkę
i są niezerowe, chyba że wskazano inaczej. Warunki łańcucha
rozważamy również dla pierścienia zerowego, w którym zachodzą
trywialnie.

\begin{definicja}[Pierścień artinowski]
\label{aa:def:artinian-ring}
Przemienny pierścień $R$ z jedynką nazywamy
\emph{artinowskim}, jeśli każdy
zstępujący ciąg jego ideałów
$I_1\supseteq I_2\supseteq\cdots$ jest od pewnego miejsca
stały, to znaczy istnieje $N$ takie, że
$I_n=I_N$ dla wszystkich $n\ge N$.
\end{definicja}

\begin{przyklad}[Pierścienie skończone, ciała i liczby całkowite]
\label{aa:prz:artin-basic}
Każdy pierścień skończony jest artinowski: w ściśle
zstępującym ciągu podzbiorów skończonego zbioru
liczności musiałyby ściśle maleć bez końca.
Dowolne ciało $K$ także jest artinowskie, nawet gdy
jest nieskończone, ponieważ ma tylko ideały $(0)$ i $K$.
Natomiast $\mathbb Z$ nie jest artinowski:
\[
 (2)\supsetneq(2^2)\supsetneq(2^3)\supsetneq\cdots.
\]
Każde zawieranie jest ścisłe, bo $2^n$ nie jest
wielokrotnością $2^{n+1}$ w $\mathbb Z$.
Analogicznie $K[x]$ nie jest artinowski dzięki
łańcuchowi $(x)\supsetneq(x^2)\supsetneq\cdots$.
Oba pierścienie są jednak noetherowskie.
\end{przyklad}

\begin{lemat}[Ilorazy pierścieni artinowskich]
\label{aa:lem:artin-quotient}
Niech $R$ będzie przemiennym pierścieniem artinowskim
z jedynką, a $I\triangleleft R$ jego ideałem.
Wówczas $R/I$ jest artinowski, również gdy $I=R$.
\end{lemat}
\begin{proof}
Niech $J_1\supseteq J_2\supseteq\cdots$ będzie
zstępującym ciągiem ideałów $R/I$.
Przeciwobrazy $\pi^{-1}(J_n)$ przez kanoniczny
homomorfizm $\pi:R\to R/I$ tworzą zstępujący ciąg
ideałów $R$, więc od pewnego miejsca są równe.
Ponieważ $\pi$ jest surjekcją, równość przeciwobrazów
implikuje równość odpowiadających im ideałów $J_n$.
\end{proof}

\section{Dziedziny artinowskie są ciałami}
\begin{twierdzenie}[Dziedzina całkowitości artinowska jest ciałem]
\label{aa:tw:artin-domain-field}
Każda dziedzina całkowitości $D$, która jest
pierścieniem artinowskim, jest ciałem.
\end{twierdzenie}
\begin{proof}
Niech $a\in D$ będzie niezerowy. Rozważmy ciąg
ideałów $(a)\supseteq(a^2)\supseteq(a^3)\supseteq\cdots$.
Z artinowskości istnieje $n\ge1$ takie, że
$(a^n)=(a^{n+1})$. W szczególności $a^n$ należy
do ideału $(a^{n+1})$, zatem dla pewnego $b\in D$
mamy $a^n=a^{n+1}b$. Przepisujemy tę równość jako
$a^n(1-ab)=0$. Ponieważ $D$ jest dziedziną i
$a^n\ne0$, możemy skrócić $a^n$, otrzymując
$ab=1$. Każdy niezerowy element $D$ ma więc
odwrotność; przemienność daje ciało.
\end{proof}

\begin{wniosek}[Ideały pierwsze pierścienia artinowskiego]
\label{aa:wn:artin-primes-maximal}
Niech $R$ będzie przemiennym pierścieniem artinowskim
z jedynką, a $P$ jego ideałem pierwszym.
Wówczas $P$ jest ideałem maksymalnym.
\end{wniosek}
\begin{proof}
Dla ideału pierwszego $P$ iloraz $R/P$ jest
dziedziną całkowitości. Na mocy
Lematu~\ref{aa:lem:artin-quotient} jest też artinowski,
więc poprzednie twierdzenie pokazuje, że jest ciałem.
Ideał $P$ jest maksymalny wtedy i tylko wtedy, gdy
iloraz $R/P$ jest ciałem.
\end{proof}

\section{Pierścienie lokalne i radykał Jacobsona}
\begin{definicja}[Radykał Jacobsona]
\label{aa:def:jacobson-radical}
Niech $R$ będzie niezerowym przemiennym pierścieniem
z jedynką. \emph{Radykałem Jacobsona} nazywamy ideał
\[
 J(R)=\bigcap_{\mathfrak m\in\operatorname{Max}(R)}\mathfrak m,
\]
gdzie $\operatorname{Max}(R)$ oznacza rodzinę wszystkich
ideałów maksymalnych pierścienia $R$.
Każdy niezerowy pierścień przemienny z jedynką ma co najmniej
jeden ideał maksymalny: stosujemy lemat Kuratowskiego--Zorna
do rodziny ideałów właściwych uporządkowanych przez zawieranie.
Rodzina jest niepusta, bo zawiera ideał zerowy. Dla
łańcucha pustego jego ograniczeniem górnym jest ideał
zerowy. Suma mnogościowa niepustego łańcucha ideałów
właściwych jest ideałem: dla dwóch jej elementów
wybieramy z łańcucha ideał zawierający oba. Pozostaje
on właściwy, gdyż obecność $1_R$ w tej sumie oznaczałaby,
że $1_R$ należy do jednego z ideałów łańcucha,
wbrew jego właściwości.
\end{definicja}

\begin{twierdzenie}[Kryterium jednostek dla radykału Jacobsona]
\label{aa:tw:jacobson-unit-criterion}
Niech $R$ będzie niezerowym przemiennym pierścieniem
z jedynką i $a\in R$. Następujące warunki są równoważne:
\begin{enumerate}[label=\textup{(\roman*)}]
\item $a\in J(R)$;
\item dla każdego $b\in R$ element $1-ab$ jest jednostką $R$.
\end{enumerate}
\end{twierdzenie}
\begin{proof}
Załóżmy $a\in J(R)$ i wybierzmy $b\in R$.
Gdyby $1-ab$ nie był jednostką, ideał $(1-ab)$
byłby właściwy i zawierałby się w pewnym ideale
maksymalnym $\mathfrak m$. Ale $a\in\mathfrak m$,
a więc także $ab\in\mathfrak m$; stąd
$1=(1-ab)+ab\in\mathfrak m$, sprzeczność.
Oznacza to, że $1-ab$ jest jednostką.

Odwrotnie, przypuśćmy, że $a\notin J(R)$.
Istnieje wtedy ideał maksymalny $\mathfrak m$ taki,
że $a\notin\mathfrak m$. Ponieważ ideał
$\mathfrak m+(a)$ ściśle zawiera $\mathfrak m$,
maksymalność daje $\mathfrak m+(a)=R$.
Istnieją zatem $c\in\mathfrak m$ i $b\in R$
takie, że $1=c+ab$. Wtedy $1-ab=c\in\mathfrak m$
i nie jest jednostką. Warunek (ii) zawodzi.
\end{proof}

\begin{twierdzenie}[Charakterystyka pierścieni lokalnych]
\label{aa:tw:local-nonunits}
Dla niezerowego przemiennego pierścienia $R$ z jedynką
następujące warunki są równoważne:
\begin{enumerate}[label=\textup{(\roman*)}]
\item $R$ ma dokładnie jeden ideał maksymalny $\mathfrak m$;
\item zbiór wszystkich elementów nieodwracalnych $R$
      jest ideałem właściwym.
\end{enumerate}
Jeżeli te warunki zachodzą, to element $a\in R$ jest
jednostką wtedy i tylko wtedy, gdy $a\notin\mathfrak m$.
\end{twierdzenie}
\begin{proof}
Załóżmy (i). Żadna jednostka nie należy do ideału
właściwego, więc każdy element $\mathfrak m$ jest
nieodwracalny. Jeśli $a\notin\mathfrak m$ byłby
nieodwracalny, ideał $(a)$ byłby właściwy i zawierałby
się w pewnym ideale maksymalnym. Jedynym takim ideałem
jest $\mathfrak m$, co przeczyłoby $a\notin\mathfrak m$.
Zbiór elementów nieodwracalnych jest więc dokładnie
$\mathfrak m$, czyli ideałem właściwym.

Załóżmy (ii) i oznaczmy ten ideał przez $N$.
Każdy ideał właściwy nie zawiera jednostek, więc
zawiera się w $N$. W szczególności $N$ jest maksymalny
i każdy inny ideał maksymalny musi być równy $N$.
\end{proof}

\begin{przyklad}[Lokalna algebra z elementami nilpotentnymi]
\label{aa:prz:local-artin-examples}
Dla ciała $K$ i liczby całkowitej $n\ge2$ pierścień
$R=K[t]/(t^n)$ jest lokalny. Element ma postać
$a_0+a_1t+\cdots+a_{n-1}t^{n-1}+(t^n)$.
Jeśli $a_0\ne0$, to jest jednostką:
po wyłączeniu $a_0$ piszemy go jako $a_0(1+u)$,
gdzie $u$ jest wielokrotnością klasy $t$ i $u^n=0$.
Odwrotnością $1+u$ jest skończona suma
$1-u+u^2-\cdots+(-u)^{n-1}$. Jeżeli $a_0=0$,
element jest nilpotentny, a zatem nie może być
jednostką w niezerowym pierścieniu. Zbiór
nieodwracalnych elementów jest więc ideałem $(t)$.
Jest to jedyny ideał maksymalny.

Pierścień $R$ jest także artinowski. Istotnie,
ideały $R$ odpowiadają ideałom $K[t]$ zawierającym
$(t^n)$. Ponieważ $K[t]$ jest dziedziną ideałów
głównych, każdy taki ideał ma postać $(f)$.
Zawieranie $(t^n)\subseteq(f)$ oznacza, że
$f$ dzieli $t^n$. Wielomian $t$ jest nierozkładalny
w $K[t]$, zatem jednoznaczność rozkładu na czynniki
nierozkładalne daje $f=ct^j$ dla pewnego $c\in K^\times$
i $0\le j\le n$. Otrzymujemy dokładnie ideały
$(t^j)$ (włącznie z całym pierścieniem i ideałem
zerowym po przejściu do ilorazu).
Mamy więc tylko skończenie wiele ideałów.
\end{przyklad}

\section{Skończoność zbioru ideałów maksymalnych}
\begin{twierdzenie}[Skończenie wiele ideałów maksymalnych]
\label{aa:tw:artin-finite-maximals}
Niezerowy przemienny pierścień artinowski $R$ ma
skończenie wiele ideałów maksymalnych.
\end{twierdzenie}
\begin{proof}
Załóżmy przeciwnie, że istnieje nieskończenie wiele
różnych ideałów maksymalnych
$\mathfrak m_1,\mathfrak m_2,\ldots$.
Połóżmy $I_n=\bigcap_{j=1}^n\mathfrak m_j$.
Jest to zstępujący ciąg ideałów. Pokażemy, że każde
zawieranie $I_{n+1}\subseteq I_n$ jest ścisłe.
Dla każdego $j\le n$ mamy
$\mathfrak m_j\not\subseteq\mathfrak m_{n+1}$:
gdyby zachodziło zawieranie, maksymalność wymuszałaby
równość tych ideałów. Wybierzmy więc
$a_j\in\mathfrak m_j\setminus\mathfrak m_{n+1}$.
Iloczyn $a_1\cdots a_n$ należy do wszystkich
$\mathfrak m_j$ dla $j\le n$, a więc do $I_n$.
Nie należy jednak do $\mathfrak m_{n+1}$, ponieważ
każdy ideał maksymalny w pierścieniu przemiennym
jest pierwszy. Stąd
$a_1\cdots a_n\in I_n\setminus I_{n+1}$.
Otrzymaliśmy nieskończony ściśle zstępujący ciąg
ideałów, sprzeczny z artinowskością $R$.
\end{proof}

\section{Nilpotentność radykału Jacobsona}
\begin{twierdzenie}[Nilpotentność radykału Jacobsona]
\label{aa:tw:artin-jacobson-nilpotent}
Niech $R$ będzie niezerowym przemiennym pierścieniem
artinowskim z jedynką. Wówczas istnieje liczba
całkowita $N\ge1$ taka, że $J(R)^N=(0)$.
\end{twierdzenie}
\begin{proof}
Piszmy $J=J(R)$. Zstępujący ciąg
$J\supseteq J^2\supseteq\cdots$ stabilizuje się,
więc dla pewnego $N\ge1$ mamy
$J^N=J^{N+1}$. Mnożąc tę równość przez kolejne
potęgi $J$, otrzymujemy $J^N=J^{N+k}$ dla każdego
$k\ge0$. Połóżmy $I=J^N$. W szczególności
$I^2=J^{2N}=J^N=I$.

Załóżmy, że $I\ne(0)$. Rozważmy niepustą rodzinę
$\mathcal F$ wszystkich ideałów $K\subseteq I$
takich, że $IK\ne(0)$. Rodzina jest niepusta, ponieważ
$I\in\mathcal F$ na mocy $I^2=I\ne(0)$.
Artinowskość zapewnia, że $\mathcal F$ ma element
minimalny $K$: gdyby żaden element nie był minimalny,
można byłoby zbudować nieskończony ściśle zstępujący
ciąg ideałów w tej rodzinie.

Skoro $IK\ne0$, istnieje $a\in K$ z $Ia\ne0$.
Ideał $(a)$ należy do $\mathcal F$, ponieważ
$I(a)=Ia\ne0$, i zawiera się w $K$, zatem
minimalność daje $K=(a)$.
Ponieważ $IK\subseteq K$, a
$I(IK)=I^2K=IK\ne0$, także $IK$ należy do
$\mathcal F$. Z minimalności $K$ wynika $IK=K$.
Stąd $a\in K=I(a)$, czyli $a=ba$ dla pewnego
$b\in I$. Ponieważ $I=J^N\subseteq J$, mamy
$b\in J$. Z
Twierdzenia~\ref{aa:tw:jacobson-unit-criterion}
wynika, że $1-b$ jest jednostką.
Równość $(1-b)a=0$ wymusza więc $a=0$, co przeczy
$Ia\ne0$. Zatem $I=(0)$, czyli $J^N=(0)$.
\end{proof}

\begin{wniosek}[Iloraz przez radykał Jacobsona]
\label{aa:wn:artin-semisimple-quotient}
Niech $R$ będzie niezerowym przemiennym pierścieniem
artinowskim, a $\mathfrak m_1,\ldots,\mathfrak m_r$
jego wszystkimi ideałami maksymalnymi. Wtedy
\[
 R/J(R)\cong\prod_{i=1}^r R/\mathfrak m_i,
\]
a każdy czynnik po prawej jest ciałem.
\end{wniosek}
\begin{proof}
Różne ideały maksymalne są parami względnie pierwsze:
$\mathfrak m_i+\mathfrak m_j=R$ dla $i\ne j$, ponieważ
suma jest ideałem ściśle zawierającym każdy z nich.
Chińskie twierdzenie o resztach daje surjekcję
$R\to\prod_i R/\mathfrak m_i$ z jądrem
$\bigcap_i\mathfrak m_i=J(R)$. Pierwsze twierdzenie
o izomorfizmie pierścieni kończy dowód.
\end{proof}

\section{Struktura przemiennych pierścieni artinowskich}
\begin{twierdzenie}[Rozkład na pierścienie artinowskie lokalne]
\label{aa:tw:artin-local-product}
Niech $R$ będzie niezerowym przemiennym pierścieniem
artinowskim z jedynką, a
$\mathfrak m_1,\ldots,\mathfrak m_r$ jego wszystkimi
ideałami maksymalnymi. Wtedy dla pewnego $N\ge1$
zachodzi izomorfizm pierścieni
\[
 R\cong\prod_{i=1}^r R/\mathfrak m_i^N.
\]
Każdy pierścień $R/\mathfrak m_i^N$ jest artinowski
i lokalny, a jego jedyny ideał maksymalny jest obrazem
$\mathfrak m_i$.
\end{twierdzenie}
\begin{proof}
Na mocy poprzednich twierdzeń ideałów maksymalnych
jest skończenie wiele, a $J=\bigcap_i\mathfrak m_i$
jest nilpotentny, powiedzmy $J^N=0$.
Przypomnijmy, dlaczego dla ideałów względnie pierwszych
$I+K=R$ zachodzi $I\cap K=IK$. Wybierzmy $u\in I$
i $v\in K$ z $u+v=1$. Jeśli $z\in I\cap K$, to
$z=zu+zv\in IK$, a odwrotne zawieranie jest oczywiste.
Ponadto, jeśli ideał $L$ jest względnie pierwszy z każdym
z ideałów $I_1,\ldots,I_q$, to jest względnie pierwszy
z ich iloczynem: z równości $u_j+v_j=1$, gdzie
$u_j\in I_j$ i $v_j\in L$, po pomnożeniu wszystkich
tych równości otrzymujemy
$1=\prod_j u_j+w$ dla pewnego $w\in L$.
Indukcja daje więc równość przecięcia i iloczynu
dla skończonej rodziny ideałów parami względnie pierwszych.
W szczególności $J=\prod_i\mathfrak m_i$.
Wobec tego
$\prod_i\mathfrak m_i^N=J^N=(0)$.

Potęgi różnych ideałów maksymalnych pozostają
względnie pierwsze. Jeśli $a\in\mathfrak m_i$,
$b\in\mathfrak m_j$ i $a+b=1$, to w rozwinięciu
$(a+b)^{2N-1}=1$ każdy składnik zawiera albo
$a^N$, albo $b^N$. Stąd
$1\in\mathfrak m_i^N+\mathfrak m_j^N$.
Chińskie twierdzenie o resztach daje więc
\[
 R/\bigcap_i\mathfrak m_i^N
 \cong\prod_{i=1}^r R/\mathfrak m_i^N.
\]
Ponieważ ideały po lewej są parami względnie pierwsze,
przecięcie jest ich iloczynem, czyli $(0)$.
Otrzymujemy żądany izomorfizm.

Każdy czynnik jest artinowski jako iloraz $R$.
Aby pokazać lokalność czynnika dla ustalonego $i$,
rozważmy jego dowolny ideał maksymalny. Jest obrazem
pewnego ideału maksymalnego $\mathfrak m_j$ pierścienia
$R$ zawierającego $\mathfrak m_i^N$.
Ponieważ $\mathfrak m_j$ jest pierwszy i
$\mathfrak m_i^N\subseteq\mathfrak m_j$, dostajemy
$\mathfrak m_i\subseteq\mathfrak m_j$:
dla $a\in\mathfrak m_i$ mamy $a^N\in\mathfrak m_j$,
a więc $a\in\mathfrak m_j$. Maksymalność daje
$\mathfrak m_i=\mathfrak m_j$.
Tylko ten ideał maksymalny występuje w ilorazie,
co dowodzi lokalności.
\end{proof}

\begin{przyklad}[Rozkład pierścienia $\mathbb Z/12\mathbb Z$]
\label{aa:prz:artin-z12}
Pierścień $\mathbb Z/12\mathbb Z$ jest artinowski,
gdyż ma skończenie wiele elementów. Jego ideały
odpowiadają ideałom $\mathbb Z$ zawierającym $(12)$,
czyli ideałom $(d)$ dla dodatnich dzielników $d\mid12$.
Redukcja modulo $d$ wyznacza surjektywny homomorfizm
$\mathbb Z/12\mathbb Z\to\mathbb Z/d\mathbb Z$
o jądrze $(\overline d)$. Stąd iloraz przez
$(\overline d)$ jest izomorficzny z
$\mathbb Z/d\mathbb Z$, a więc jest ciałem dokładnie
wtedy, gdy $d$ jest liczbą pierwszą. Stąd mamy dokładnie
dwa ideały maksymalne: $(\overline2)$ i $(\overline3)$.
Ich przecięcie to ideał $(\overline6)$, składający
się z klas $\overline0,\overline6$.
Jest to radykał Jacobsona, a jego kwadrat wynosi
$(\overline{36})=(\overline0)$. Chińskie twierdzenie o resztach daje izomorfizm
$\mathbb Z/12\mathbb Z\cong
\mathbb Z/4\mathbb Z\times\mathbb Z/3\mathbb Z$,
$\overline a\mapsto(a\bmod4,a\bmod3)$.
Jego odwrotność wysyła parę $(u,v)$ na klasę
$9u+4v$ modulo $12$, ponieważ $9\equiv1\pmod4$,
$9\equiv0\pmod3$, $4\equiv0\pmod4$ i
$4\equiv1\pmod3$.
W opisie z Twierdzenia~\ref{aa:tw:artin-local-product}
możemy przyjąć $N=2$: dla ideałów maksymalnych
$\mathfrak m_2=(\overline2)$ i $\mathfrak m_3=(\overline3)$
mamy $\mathfrak m_2^2=(\overline4)$ oraz
$\mathfrak m_3^2=(\overline9)=(\overline3)$.
Dlatego czynniki $R/\mathfrak m_2^2$ i
$R/\mathfrak m_3^2$ są izomorficzne odpowiednio
z $\mathbb Z/4\mathbb Z$ i $\mathbb Z/3\mathbb Z$.
Pierwszy czynnik, $\mathbb Z/4\mathbb Z$, jest lokalny
z ideałem maksymalnym generowanym przez klasę $2$
modulo $4$, a drugi jest ciałem.
Ten przykład pokazuje, jak nilpotentność radykału
współistnieje z rozkładem pierścienia na lokalne czynniki.
\end{przyklad}

\section{Dlaczego pierścień artinowski jest noetherowski?}
\begin{lemat}[Moduły artinowskie nad skończonym iloczynem ciał]
\label{aa:lem:artin-mod-over-fields}
Niech $r\ge1$ będzie liczbą całkowitą, a
$A=K_1\times\cdots\times K_r$ skończonym iloczynem ciał i niech $M$ będzie $A$-modułem
spełniającym warunek stabilizacji zstępujących
ciągów podmodułów. Wtedy $M$ jest skończenie generowany
jako $A$-moduł.
\end{lemat}
\begin{proof}
Niech $e_i\in A$ będzie elementem równym $1$
na $i$-tej współrzędnej i $0$ na pozostałych.
Wtedy $e_i^2=e_i$, $e_ie_j=0$ dla $i\ne j$
oraz $\sum_i e_i=1$. Każdy $m\in M$ ma więc
jednoznaczny rozkład $m=\sum_i e_im$, a
$M=\bigoplus_i e_iM$.
Moduł $e_iM$ jest przestrzenią wektorową nad $K_i$:
działanie $A$ na nim zależy wyłącznie od $i$-tej
współrzędnej skalara. Każda jego podprzestrzeń $W\subseteq e_iM$ jest też
$A$-podmodułem $M$: dla $a=(a_1,\ldots,a_r)\in A$
i $w\in W$ zachodzi $aw=a_iw\in W$. Dlatego
każdy zstępujący ciąg takich podprzestrzeni stabilizuje się.
Gdyby $e_iM$ było nieskończenie wymiarowe,
można byłoby wybrać liniowo niezależne wektory
$v_1,v_2,\ldots$. Podprzestrzenie
$\operatorname{span}\{v_n,v_{n+1},\ldots\}$
tworzyłyby nieskończony ściśle zstępujący ciąg,
co jest sprzecznością. Zatem każdy $e_iM$ ma
skończony wymiar nad $K_i$, a złączenie skończonych
baz tych przestrzeni generuje $M$ jako $A$-moduł.
\end{proof}

\begin{twierdzenie}[Pierścień artinowski jest noetherowski]
\label{aa:tw:artin-implies-noether}
Każdy przemienny pierścień artinowski z jedynką
jest noetherowski.
\end{twierdzenie}
\begin{proof}
Przypadek pierścienia zerowego jest natychmiastowy;
załóżmy więc, że $R\ne0$. Niech $J=J(R)$.
Na mocy Twierdzenia~\ref{aa:tw:artin-jacobson-nilpotent}
przy pewnym $N$ zachodzi $J^N=0$.
Z Wniosku~\ref{aa:wn:artin-semisimple-quotient}
wiemy, że $A=R/J$ jest skończonym iloczynem ciał.

Ustalmy dowolny ideał $I\subseteq R$.
Rozważmy skończoną filtrację
$ I=I\cap J^0\supseteq I\cap J\supseteq\cdots
\supseteq I\cap J^N=0$.
Dla każdego $k=0,\ldots,N-1$ moduł
\[
 M_k=(I\cap J^k)/(I\cap J^{k+1})
\]
można utożsamić z podmodułem $J^k/J^{k+1}$:
odwzorowanie $a+(I\cap J^{k+1})\mapsto a+J^{k+1}$
jest dobrze określone i ma jądro zerowe, gdyż
$a\in I\cap J^k$ oraz $a\in J^{k+1}$ oznaczają
$a\in I\cap J^{k+1}$.
Mnożenie przez element $J$ zeruje każdą klasę w $M_k$,
ponieważ $J(I\cap J^k)\subseteq I\cap J^{k+1}$.
Działanie $R$ przechodzi więc na $A=R/J$.

Moduł $M_k$ spełnia warunek stabilizacji zstępujących
ciągów $A$-podmodułów. Każdy taki podmoduł jest też
$R$-podmodułem, ponieważ działanie $R$ odbywa się
poprzez homomorfizm $R\to A$. Jego przeciwobraz przy
rzutowaniu $I\cap J^k\to M_k$ jest zatem ideałem
pierścienia $R$ zawierającym $I\cap J^{k+1}$.
Zstępujące ciągi tych przeciwobrazów stabilizują się,
bo $R$ jest artinowski. Surjektywność rzutowania
zapewnia wówczas stabilizację wyjściowego ciągu.
Z Lematu~\ref{aa:lem:artin-mod-over-fields} wynika,
że każdy $M_k$ ma skończenie wielu generatorów.

Wybierzmy dla każdego $k$ skończenie wiele reprezentantów
w $I\cap J^k$ generatorów $M_k$. Ich łączna rodzina
jest skończona i generuje $I$ jako ideał: dla danego
$a\in I$ wyrażamy jego klasę w $M_0$ jako kombinację
obrazów reprezentantów ze współczynnikami w $A$,
a współczynniki podnosimy do $R$. Po odjęciu tej
kombinacji pozostaje element $I\cap J$.
Powtarzając to samo dla jego klasy w $M_1$, otrzymujemy
element $I\cap J^2$, i tak dalej.
Po co najwyżej $N$ krokach reszta należy do
$I\cap J^N=0$. Wszystkie odejmowane składniki należą
do ideału generowanego przez wybrane reprezentanty.
Zatem każdy ideał $I$ jest skończenie generowany.
\end{proof}

\section{Lemat Nakayamy i minimalne układy generatorów}
\begin{twierdzenie}[Lemat Nakayamy dla pierścienia lokalnego]
\label{aa:tw:nakayama-local}
Niech $(R,\mathfrak m)$ będzie przemiennym pierścieniem
lokalnym z jedynką i jedynym ideałem maksymalnym
$\mathfrak m$. Niech $M$ będzie skończenie generowanym
modułem nad $R$. Jeżeli $\mathfrak m M=M$,
to $M=0$.
\end{twierdzenie}
\begin{proof}
Załóżmy, że $M\ne0$. Spośród wszystkich skończonych
układów generatorów $M$ wybierzmy taki, który ma
najmniejszą liczbę $r\ge1$ elementów,
powiedzmy $m_1,\ldots,m_r$.
Ponieważ $M=\mathfrak m M$, element $m_1$
można zapisać jako
$m_1=a_1m_1+\cdots+a_rm_r$ z $a_i\in\mathfrak m$.
Stąd
$(1-a_1)m_1=a_2m_2+\cdots+a_rm_r$.

Element $1-a_1$ nie należy do $\mathfrak m$:
gdyby należał, to $1=(1-a_1)+a_1\in\mathfrak m$.
Na mocy Twierdzenia~\ref{aa:tw:local-nonunits}
element $1-a_1$ jest więc jednostką.
Mnożąc ostatnią równość przez jego odwrotność,
otrzymujemy, że $m_1$ należy do podmodułu
generowanego przez $m_2,\ldots,m_r$.
Wtedy te pozostałe elementy generują całe $M$,
co przeczy minimalności $r$.
Wobec tego $M=0$.
\end{proof}

\begin{wniosek}[Wersja Nakayamy z podmodułem]
\label{aa:wn:nakayama-submodule}
Niech $(R,\mathfrak m)$ będzie niezerowym przemiennym
pierścieniem lokalnym z jedynką i jedynym ideałem
maksymalnym $\mathfrak m$, niech $M$ będzie skończenie
generowanym $R$-modułem, a
$N\subseteq M$ podmodułem takim, że
$M=N+\mathfrak m M$. Wtedy $M=N$.
\end{wniosek}
\begin{proof}
Moduł ilorazowy $M/N$ jest skończenie generowany,
bo obrazy dowolnych skończonych generatorów $M$
generują $M/N$. Założenie
$M=N+\mathfrak m M$ po przejściu do ilorazu daje
$M/N=\mathfrak m(M/N)$. Twierdzenie Nakayamy
zastosowane do $M/N$ prowadzi do $M/N=0$.
Jest to równoważne $M=N$.
\end{proof}

\begin{twierdzenie}[Minimalna liczba generatorów modułu]
\label{aa:tw:nakayama-minimal-generators}
Niech $(R,\mathfrak m)$ będzie przemiennym pierścieniem
lokalnym z jedynką, $K=R/\mathfrak m$ jego ciałem
reszt, a $M$ skończenie generowanym $R$-modułem.
Wtedy $M/\mathfrak m M$ jest skończenie wymiarową
przestrzenią wektorową nad $K$, a najmniejsza możliwa
liczba generatorów $M$ jest równa
$\dim_K(M/\mathfrak m M)$.
\end{twierdzenie}
\begin{proof}
Ponieważ $\mathfrak m$ działa zerowo na
$M/\mathfrak m M$, mnożenie przez $R$ zależy tylko
od klasy skalara modulo $\mathfrak m$.
W ten sposób iloraz jest przestrzenią wektorową nad $K$.
Obrazy skończonej rodziny generatorów $M$ rozpinają
ten iloraz, więc jego wymiar jest skończony, powiedzmy $d$.

Jeżeli $M$ ma $r$ generatorów, ich obrazy rozpinają
przestrzeń wymiaru $d$, zatem $r\ge d$.
Aby wykazać, że $d$ generatorów wystarczy,
wybierzmy bazę $\bar m_1,\ldots,\bar m_d$
przestrzeni $M/\mathfrak m M$ i jej reprezentantów
$m_i\in M$. Niech $N=(m_1,\ldots,m_d)$.
Każda klasa modulo $\mathfrak m M$ jest kombinacją
$\bar m_i$, więc $M=N+\mathfrak m M$.
Z Wniosku~\ref{aa:wn:nakayama-submodule} otrzymujemy
$M=N$. Zatem $d$ jest zarazem osiągalną i najmniejszą
liczbą generatorów.
\end{proof}

\begin{przyklad}[Różnica między bazą a generatorami modułu]
\label{aa:prz:nakayama-example}
Niech $K$ będzie ciałem,
$R=K[x]_{(x)}$ lokalizacją pierścienia wielomianów
w jego ideale maksymalnym $(x)$. Jest to pierścień
lokalny z ideałem maksymalnym $\mathfrak m=(x)R$.
Dla modułu wolnego $M=R^2$ mamy
$M/\mathfrak m M\cong K^2$.
Poprzednie twierdzenie mówi, że $M$ wymaga
co najmniej dwóch generatorów i standardowe dwa
wektory bazowe rzeczywiście wystarczają.

Z kolei dla $M=R/(x^3)R$ iloraz
$M/\mathfrak m M$ jest izomorficzny z $K$,
więc wystarczy jeden generator: klasa $1$.
Nie jest to jednak wolny moduł rangi jeden:
niezerowy element $x^3\in R$ anihiluje cały moduł
$M$, bo $x^3(q+(x^3)R)=0$ dla każdego $q\in R$.
Natomiast w wolnym module rangi jeden, izomorficznym
z $R$, żaden niezerowy element $r\in R$ nie anihiluje
całego modułu, ponieważ $r\cdot1=r\ne0$.
Należy odróżniać \emph{posiadanie
jednego generatora} od \emph{posiadania bazy złożonej
z jednego elementu}.
\end{przyklad}


\begin{twierdzenie}[Lemat Nakayamy dla radykału Jacobsona]
\label{aa:tw:nakayama-jacobson-general}
Niech $R$ będzie niezerowym przemiennym pierścieniem
z jedynką, $J(R)$ jego radykałem Jacobsona, a $M$ skończenie
generowanym $R$-modułem. Jeśli ideał
$I\subseteq J(R)$ spełnia $IM=M$, to $M=0$.
\end{twierdzenie}
\begin{proof}
Przypuśćmy, że $M\neq0$, i wybierzmy układ
generatorów $m_1,\ldots,m_r$ o minimalnej
liczbie elementów $r\ge1$. Ponieważ $IM=M$,
istnieją $a_1,\ldots,a_r\in I$ takie, że
$m_1=\sum_{i=1}^r a_im_i$.
Po przeniesieniu pierwszego składnika otrzymujemy
$(1-a_1)m_1=\sum_{i=2}^r a_im_i$.
Ponieważ $a_1\in I\subseteq J(R)$, z
Twierdzenia~\ref{aa:tw:jacobson-unit-criterion}
(dla $b=1$) wynika, że $1-a_1$ jest jednostką.
Dlatego $m_1$ jest kombinacją $m_2,\ldots,m_r$,
a cały moduł $M$ byłby generowany przez $r-1$
elementów. Jest to sprzeczne z minimalnością $r$.
Zatem $M=0$.
\end{proof}

\begin{przyklad}[Skończone generowanie jest niezbędne]
\label{aa:prz:nakayama-non-fg-counterexample}
Niech $K$ będzie ciałem,
$R=K[x]_{(x)}$ lokalnym pierścieniem funkcji
wymiernych regularnych w zerze,
a $F=K(x)$ jego ciałem ułamków.
Rozważmy $R$-moduł $M=F/R$.
Jest on niezerowy: klasa $1/x+R$ jest niezerowa,
ponieważ $x$ nie ma odwrotności w $R$.
Ideałem maksymalnym $R$ jest $\mathfrak m=xR$.
Dla każdego $q+R\in M$ mamy
$q+R=x(q/x+R)$, więc $\mathfrak m M=M$.
Mimo to $M\ne0$, co nie przeczy lematowi Nakayamy,
gdyż $M$ nie jest skończenie generowany.

Zweryfikujmy ostatnie stwierdzenie bez odwoływania
się do lematu. Każdy element $F$ jest postaci
$r/x^n$ dla pewnego $r\in R$ i $n\ge0$:
po wyłączeniu największej potęgi $x$ z mianownika
pozostały mianownik jest jednostką $R$.
Gdyby klasy $q_1+R,\ldots,q_s+R$ generowały $M$,
można byłoby dobrać wspólne $N\ge0$ tak,
aby $x^Nq_i\in R$ dla wszystkich $i$.
Wówczas $x^NM=0$, ponieważ mnożenie przez $x^N$
zerowałoby każdy generator i wszystkie jego
$R$-kombinacje. Jednak
$x^N(1/x^{N+1}+R)=1/x+R\ne0$.
Otrzymana sprzeczność dowodzi, że $M$ nie ma
skończonego układu generatorów.
\end{przyklad}

\begin{przyklad}[Dwuwymiarowa przestrzeń generatorów ideału]
\label{aa:prz:nakayama-two-generators-square-zero}
Niech $K$ będzie ciałem i
$R=K[x,y]/(x,y)^2$; oznaczmy klasy $x,y$
przez $\bar x,\bar y$. Każda klasa $R$ ma
jednoznaczną postać $a+b\bar x+c\bar y$
z $a,b,c\in K$, ponieważ jednomiany łącznego stopnia
co najmniej dwa znikają w ilorazie.
Element o $a\ne0$ jest jednostką: jeżeli
$u=b\bar x+c\bar y$, to $u^2=0$ i
$(a+u)^{-1}=a^{-1}-a^{-2}u$.
Elementy z $a=0$ tworzą ideał
$\mathfrak m=(\bar x,\bar y)$, a
$\mathfrak m^2=0$. Są więc nilpotentne i nieodwracalne;
razem z wcześniejszą obserwacją oznacza to, że zbiór
elementów nieodwracalnych jest dokładnie ideałem
$\mathfrak m$. Z Twierdzenia~\ref{aa:tw:local-nonunits}
wynika, że $R$ jest lokalny, a $\mathfrak m$ jest
jego jedynym ideałem maksymalnym.

Traktujmy teraz $M=\mathfrak m$ jako $R$-moduł.
Wówczas $\mathfrak m M=\mathfrak m^2=0$, stąd
$M/\mathfrak m M=\mathfrak m$.
Jako przestrzeń wektorowa nad
$R/\mathfrak m\cong K$ ma ona bazę
$\bar x,\bar y$ i wymiar dwa.
Twierdzenie~\ref{aa:tw:nakayama-minimal-generators}
wykazuje, że ideał $\mathfrak m$ wymaga
\emph{dokładnie dwóch} generatorów.
Jeden element $a\bar x+b\bar y$ nie może go
generować, ponieważ jego wielokrotności przez $R$
tworzą najwyżej jednowymiarową podprzestrzeń
$K\bar x+K\bar y$: składniki z $\mathfrak m$
zerują $\mathfrak m$.
\end{przyklad}

\begin{uwaga}[Przejście do szerszej teorii modułów]
\label{aa:uw:artin-to-modules}
Lemat Nakayamy jest pierwszym przykładem wyniku,
w którym struktura pierścienia lokalnego kontroluje
minimalną liczbę generatorów modułu.
W dalszych rozdziałach rozszerzymy tę teorię o
moduły wolne, prezentacje, rozszczepienia, iloczyn
tensorowy modułów, postać Smitha i klasyfikację modułów
nad dziedzinami ideałów głównych. Nie utożsamiamy
iloczynu tensorowego z iloczynem prostym.
\end{uwaga}

\setlength{\parskip}{4pt}
